\documentclass[a4paper,14pt]{extarticle}

% -------------------- ШРИФТЫ (XeLaTeX / LuaLaTeX) --------------------
\usepackage{fontspec}
\usepackage{polyglossia}

\setdefaultlanguage{russian}
\setotherlanguage{english}

% Включаем поддержку стандартных лигатур TeX (--- превратится в тире)
\setmainfont{Times New Roman}[Ligatures=TeX]
\newfontfamily{\russianfont}{Times New Roman}[Script=Cyrillic,Ligatures=TeX]
\setsansfont{Arial}[Ligatures=TeX]
\setmonofont{Courier New}
\newfontfamily{\cyrillicfonttt}{Courier New}

% -------------------- ГЕОМЕТРИЯ СТРАНИЦЫ --------------------
\usepackage{geometry}
\geometry{
  top=20mm,
  bottom=20mm,
  left=30mm,
  right=10mm
}

\sloppy
\usepackage{tabularx}
\usepackage{booktabs}
\usepackage{listings}
\usepackage{xcolor}

\usepackage{enumitem}
\setlist[enumerate]{label=--, leftmargin=*, align=left}

%\usepackage{pdfpages}
\usepackage{titlesec}
\usepackage{indentfirst}
\usepackage{seqsplit}
\usepackage{caption}
\usepackage{amsmath}
\usepackage[htt]{hyphenat}
\usepackage{graphicx}
\usepackage{float}
\usepackage[hidelinks]{hyperref}
\usepackage{setspace}
\usepackage{etoolbox}
\usepackage[
  backend=biber,
  style=gost-numeric,
  sorting=none,
  language=auto,
  autolang=other,
  bibencoding=utf8,
  url=true,
  doi=true,
  isbn=false,
  eprint=false
]{biblatex}
\DeclareFieldFormat{url}{\textnormal{\url{#1}}}
\addbibresource{include/references.bib}

\onehalfspacing % Полуторный интервал для основного текста
\setlength{\parindent}{1.25cm} % Абзацный отступ по ГОСТ

% Нумерация таблиц по разделам (например: 1.1, 1.2)
\counterwithin{table}{section}

% -------------------- ОФОРМЛЕНИЕ ЗАГОЛОВКОВ --------------------
\setcounter{secnumdepth}{3}
\setcounter{tocdepth}{2}

\titleformat{\section}
  {\centering\normalfont\normalsize\bfseries}
  {\thesection}{1em}{}

\titleformat{\subsection}
  {\centering\normalfont\normalsize\bfseries}
  {\thesubsection}{1em}{}

\titleformat{\subsubsection}
  {\centering\normalfont\normalsize\bfseries}
  {\thesubsubsection}{1em}{}

% -------------------- НАСТРОЙКА ТАБЛИЦ ПО ГОСТ --------------------
\usepackage{xltabular} % Заменяет tabularx + добавляет поддержку longtable
\AtBeginEnvironment{xltabular}{\singlespacing} % Одинарный интервал внутри таблицы по ГОСТ
% Оформление подписи: Таблица 1.1 — Название таблицы (без точки в конце)
\DeclareCaptionLabelFormat{gosttable}{Таблица #2}

% Объявляем собственный разделитель: " — " (пробел, длинное тире, пробел)
\DeclareCaptionLabelSeparator{gostdash}{ --- } 

\captionsetup[table]{
  labelformat=gosttable,
  labelsep=gostdash,         % Применяем созданный разделитель вместо emdash
  justification=raggedright, % Выравнивание по левому краю
  singlelinecheck=false,     % Отмена центрирования коротких строк
  font=normal,               % Прямой шрифт текста по ГОСТ
  position=top,              % Подпись сверху
  skip=1pt                   % Отступ до таблицы
}

% Автоматический сброс интервала до одинарного внутри таблиц (требование ГОСТ)
\AtBeginEnvironment{tabular}{\singlespacing}
\AtBeginEnvironment{tabularx}{\singlespacing}
\renewcommand{\arraystretch}{1.2} % Небольшой внутренний отступ для читаемости

% -------------------- ЛИСТИНГИ --------------------
% Оформление подписи листингов
\DeclareCaptionFormat{listing}{
  \parbox{\textwidth}{#1#2#3}
}
\AtBeginDocument{\counterwithin{lstlisting}{section}}
\renewcommand{\lstlistingname}{Листинг}

\captionsetup[lstlisting]{
  format=listing,
  labelfont=it,
  textfont=it,
  justification=raggedright,
  singlelinecheck=false
}

% Стиль кода
\lstset{
    language=bash,
    basicstyle=\ttfamily\small,
    breaklines=true,
    frame=single,
    backgroundcolor=\color{gray!5},
    keywordstyle=\color{blue!70!black},
    commentstyle=\color{gray},
    captionpos=t,
    showspaces=false,
    showstringspaces=false
}

% Автоматическое ограничение размеров изображений
\setkeys{Gin}{
    width=1\textwidth,
    height=0.45\textheight,
    keepaspectratio
}

% Настройка подписей
\usepackage[labelsep=endash]{caption}

% -------------------- ПРИЛОЖЕНИЯ --------------------
\newcounter{appendixctr}
\renewcommand{\theappendixctr}{\Asbuk{appendixctr}}

\newcommand{\appendixsection}[1]{
  \clearpage
  \refstepcounter{appendixctr}
  \section*{Приложение \theappendixctr. #1}
  \addcontentsline{toc}{section}{Приложение \theappendixctr. #1}
}
